% !TEX root = thesis.tex
\chapter{Related works}\label{sec:related_works}

\section{Deep Learning for Wireless Communications}

The integration of Deep Learning (DL) in the telecommunications field is a prominent research trend. For example, Deep Neural Networks (DNNs) were able to overcome the computational complexity of iterative solvers to perform power allocation \cite{sun2018learning}.

It has been demonstrated that DL can replace classical MMSE estimators in OFDM systems \cite{ye2018power}. Furthermore, researchers have developed communication systems by framing component replacement as an end-to-end reconstruction task \cite{oshea2017introduction}. Most of these developments have occurred within the last decade; thus, the application of DL in this field is nascent, leaving substantial room for innovation.

Additionally, new software libraries \cite{sionna} allow for the link-level simulation and dataset generation necessary to accelerate model development. It is now possible to produce synthetic data from real-world scenarios on demand with customized parameters, which solves one of the biggest bottlenecks in machine learning.

\section{Foundation Models for CSI}

Recently, researchers have attempted to implement Foundation Models (FMs) for wireless communications. This concept, which is common in other domains \cite{gpt3}, could accelerate the adoption of DL solutions across different tasks by lowering the data requirements and computational complexity of downstream models, effectively creating a universal feature extractor for Channel State Information (CSI).

The first two FMs introduced were the LWM \cite{alikhani2025lwm} and BERT4MIMO \cite{catak2025bert4mimofoundationmodelusing}. Both utilized a transformer-based architecture inspired by BERT \cite{devlin2018bert}. For this work, the LWM was chosen because:

\begin{itemize}
    \item It has gained more relevance and traction within the research community.
    \item It was pre-trained on a higher-volume dataset encompassing different numbers of subcarriers and antenna arrays.
    \item It was rigorously evaluated across distinct downstream tasks, including Line-of-Sight (LoS) and Non-Line-of-Sight (NLoS) classification, as well as beam selection.
\end{itemize}

However, it is worth noting that BERT4MIMO introduces several relevant innovations. The model was pre-trained using two distinct scenarios: stationary and dynamic user equipment, which makes the training samples significantly more diverse. Furthermore, the model architecture incorporates multiple time steps for the same subcarrier and supports configurations with multiple antennas at the user equipment.

Later, 6G WavesFM \cite{wavesfm} was developed, focusing primarily on sensing and localization. The model was pre-trained on a wider variety of data, including RF spectrograms, Wi-Fi CSI samples, and 5G CSI samples. It also introduced a fine-tuning framework evaluated across four downstream tasks: human activity sensing, RF signal classification, 5G New Radio (NR) positioning, and MIMO-OFDM channel estimation. Similar to the LWM, its training strategy is based on reconstructing masked inputs using patch-based pre-processing.

Other works have enhanced traditional Large Language Models (LLMs) and vision models using communications data \cite{Zhang_2026}. However, these approaches are limited because they adapt pre-trained models (such as ChatGPT) via instruction fine-tuning and retrieval-augmented generation (RAG) to solve specific telecommunications queries. Consequently, they act as domain-specific assistants rather than the universal physical-layer feature extractors that foundation models aim to be.

Additional recent models include WiFo \cite{Liu2025} and MultimodalWFM \cite{aboulfotouh2026multimodalwirelessfoundationmodels}, which were not chosen as benchmarks due to their architectural similarities with the LWM and their comparatively lower traction within the research community. Nevertheless, they introduce relevant novelties, such as multimodality supporting RF spectrograms, IQ streams, and CSI data.

\section{Autoencoder-based CSI Compression}

Another prominent research trend is the use of autoencoders to reduce the dimensionality of the channel matrix. CSINet \cite{wen2018deep} was the first model to achieve high compression ratios with minimal information loss. Later, a newer model based on transformers \cite{CsiTransformer} demonstrated that the self-attention mechanism is highly capable of performing compression.

These autoencoders were created to reduce the overhead caused by the necessity to transmit the entire channel matrix from the UE to the BS. As the complexity of massive MIMO systems increases, the size of the channel matrix also increases, which adds significant transmission overhead. In this work, we aim to analyze whether the compressed latent vector generated by an autoencoder is sufficient to successfully perform the same downstream tasks as the LWM.

\section{Summary and Research Gaps}

Despite the initial success of FMs for telecommunications, there remain several research gaps concerning these models. The practical computational cost of both models is unknown, and in time-constrained applications such as real-time transmission, this is a critical factor. Additionally, the LWM creates embeddings with a higher dimensionality than the original channel state information. For this reason, it is important to quantify whether the additional computational and memory costs translate into meaningful performance gains.

Conversely, the use of compressed vectors generated by autoencoders for these specific downstream tasks is a novelty. It is a new approach implemented in this paper that needs to be evaluated across different downstream tasks, similarly to the LWM, to verify its feasibility.

Finally, we seek to assess how well the representations created by the proposed autoencoder and the LWM perform in a new downstream task inspired by the power allocation models mentioned previously \cite{sun2018learning, eisen2019learning}.